\section{Predykaty pierwotnie rekurencyjne}
Relacja \(R \subset \natural^{k+1}\) jest \textbf{relacją regularną} wtw:
\[ \forall_{\vec{x} \in \natural^k} \ \exists_{y \in \natural} \ (\vec{x},y) \in R \]
Predykat \(R \subseteq \natural^k\) jest pierwotnie rekurencyjny, gdy istnieje pierwotnie rekurencyjna funkcja charakterystyczna \(f_R: \natural^k \to \natural\)
\[
f_R(\vec{x}) = 
\begin{cases}
    1 \text{ gdy } \vec{x} \in R \\
    0 \text{ gdy } \vec{x} \notin R
\end{cases}
\]
\begin{theorem}
    Rodzina predykatów pierwotnie rekurencyjnych jest zamknięta na operacje logiczne.
\end{theorem}
\begin{proof}
Definiujemy operacje logiczne za pomocą operacji z \(PRec\).
    \begin{itemize}
        \item\(a \land b \coloneq a \cdot b\) 
        \item\(\neg a \coloneq 1 \dotminus a\)
        \item\(a \lor b \coloneq \neg(\neg a \land \neg b)  \)
    \end{itemize}
\end{proof}

\begin{theorem}
    Rodzina predykatów pierwotnie rekurencyjnych jest zamknięta na kwantyfikatory ograniczone.
\end{theorem}
\begin{proof}
Definiujemy kwantyfikatory ograniczone używając \(\mu\)-rekursji:
\[\exists_{x \leq a} f(x) := sgn(f(0)+\mu x_{\leq a} (\neg f(x))\]
\[\forall_{x\leq a} f(x) := \overline{sgn} (\neg f(0)+\mu x_{\leq a} (f(x))\]
\end{proof}

